\begin{figure}[ht]
\centering
\begin{tikzpicture}[x=1cm, y=1cm, font=\small]

% A cubic-like curve showing concave-up, inflection, concave-down,
% with the tangent lying below (convex) and above (concave) the curve.
% f(x) = (1/6) x^3 - (1/2) x  on roughly [-3, 3], shifted up.

\def\f(#1){(0.16*(#1)*(#1)*(#1) - 0.5*(#1) + 2.0)}

% Axes
\draw[->, thick] (-3.4, 0) -- (3.6, 0) node[right] {$x$};
\draw[->, thick] (0, 0) -- (0, 4.2) node[above] {$y = f(x)$};

% Curve
\draw[thick, engblue] plot[domain=-3.2:3.2, samples=100]
  (\x, {\f(\x)});

% Inflection at x = 0 for this cubic
\fill[engblack] (0, {\f(0)}) circle (1.8pt);
\node[engblack, anchor=south east, font=\scriptsize] at (-0.1, {\f(0)+0.05}) {inflection $f''=0$};

% Concave-down region (x < 0): tangent above curve. Pick xa = -2.
\def\xa{-2.0}
\pgfmathsetmacro{\fa}{0.16*\xa*\xa*\xa - 0.5*\xa + 2.0}
\pgfmathsetmacro{\fpa}{0.48*\xa*\xa - 0.5}
\draw[thick, engred] (\xa-1.0, {\fa - \fpa}) -- (\xa+1.0, {\fa + \fpa});
\fill[engred] (\xa, {\fa}) circle (1.6pt);
\node[engred, anchor=south, font=\scriptsize] at (\xa, {\fa+0.45}) {$f''<0$ (concave)};

% Concave-up region (x > 0): tangent below curve. Pick xb = 2.
\def\xb{2.0}
\pgfmathsetmacro{\fb}{0.16*\xb*\xb*\xb - 0.5*\xb + 2.0}
\pgfmathsetmacro{\fpb}{0.48*\xb*\xb - 0.5}
\draw[thick, engamber] (\xb-1.0, {\fb - \fpb}) -- (\xb+1.0, {\fb + \fpb});
\fill[engamber] (\xb, {\fb}) circle (1.6pt);
\node[engamber, anchor=north, font=\scriptsize] at (\xb, {\fb-0.4}) {$f''>0$ (convex)};

\end{tikzpicture}
\caption[The second derivative reads concavity]{Where $f'' > 0$ the
graph curves upward and lies above its tangent line (convex); where
$f'' < 0$ the graph curves downward and lies below its tangent line
(concave). The transition point, where $f''$ changes sign, is an
inflection point. The convexity reading is the working tool of
one-variable optimisation: a critical point in a convex region is a
minimum, a critical point in a concave region is a maximum, and the
second-derivative test of section~5.6 formalises the picture.}
\label{fig:vol02:ch05:concavity}
\end{figure}
